%% kapa: exam=1 round=22 year=2011 period=1 form=A subject=civil_law subject_ko=민법 seq=30 number=30
%% source: Q-Net 감정평가사 기출문제 (cst00309) artlSeq=5036790 file=extracted/1008/5036790/2011년 제22회 감정평가사 1차시험 기출문제/1교시.hwp
%% conversion: kordoc 4.24.51 -> tools/gen1.py (auto); review=proofread-sonnet; flags=none
\begin{question}{30}
추인에 관한 설명으로 옳지 않은 것은? (다툼이 있으면 판례에 의함)
\begin{choices}
  \item 당사자가 무효임을 알고 추인한 때에는 반사회적 법률행위도 유효로 될 수 있다.
  \item 취소할 수 있는 법률행위를 추인한 후에는 취소하지 못한다.
  \item 가장매매의 당사자가 무효인 매매를 추인하면 새로운 매매계약을 체결한 것으로 본다.
  \item 법정대리인은 취소의 원인이 종료하기 전에도 취소할 수 있는 법률행위를 추인할 수 있다.
  \item 원인무효인 등기가 이루어진 사실을 알고도 장기간 이의를 한 바 없다는 사실만으로 이를 추인한 것으로는 볼 수 없다.
\end{choices}
\end{question}
\answer{1}
